\chapter{Marco teórico y trabajos relacionados}
\label{cap:marco}

\section{Modelos de riesgo clínico y su evaluación}

Un modelo de estratificación de riesgo estima la probabilidad de un desenlace
--aquí, la defunción-- a partir de información disponible en el momento en que se
toma la decisión clínica. Su evaluación exige separar dos propiedades que se
confunden con frecuencia.

La \textbf{discriminación} mide la capacidad de ordenar pacientes por riesgo y se
resume en el AUROC. Es insensible a transformaciones monótonas del puntaje: un
modelo puede discriminar perfectamente y, aun así, emitir probabilidades
absurdas. La \textbf{calibración} mide justamente eso: si de cada cien pacientes
a quienes el modelo asigna un riesgo de 20\,\% fallecen efectivamente veinte. Se
ha argumentado con razón que la calibración es el talón de Aquiles de la
analítica predictiva en medicina, porque un modelo descalibrado desplaza
sistemáticamente las decisiones de umbral sin que el AUROC lo
delate~\parencite{vancalster2019}. Esta distinción es central en el
capítulo~\ref{cap:resultados}: el modelo de referencia de esta tesis discrimina
bien desde el principio, pero requiere recalibrado explícito antes de poder
auditarse.

En cohortes con desenlaces poco frecuentes, el AUPRC complementa al AUROC porque
es sensible al desbalance de clases. Y como toda métrica agregada oculta la
distribución de errores, la evaluación de un modelo destinado al despliegue debe
reportarse en el \emph{punto de operación} previsto y no en un umbral arbitrario.
Las guías de reporte de modelos predictivos en salud recogen estos requisitos de
manera sistemática~\parencite{tripod2015}, y la revisión crítica de modelos
pronósticos de COVID-19 documentó que su incumplimiento fue masivo durante la
pandemia: la mayoría de los modelos revisados presentaban alto riesgo
de sesgo, principalmente por selección de cohorte, fuga de información y ausencia
de validación~\parencite{wynants2020}.

\section{Aprendizaje automático sobre datos tabulares}

Para datos tabulares heterogéneos, los ensambles de árboles con impulso del
gradiente siguen siendo el punto de referencia práctico frente a arquitecturas
neuronales. Las implementaciones basadas en histogramas agrupan los valores de
cada variable en cubetas antes de buscar los cortes, lo que reduce el costo de
entrenamiento de forma sustancial y las hace viables en millones de
registros~\parencite{lightgbm2017}. En esta tesis se emplea
\texttt{HistGradientBoostingClassifier} de \emph{scikit-learn}~\parencite{sklearn2011},
que además maneja valores faltantes de forma nativa. Es una propiedad relevante
aquí, porque la codificación de ausencia de la base de la DGE resulta informativa
--el capítulo~\ref{cap:resultados} lo cuantifica-- y no conviene imputarla a
ciegas.

El desbalance de clases se atiende reponderando la función de pérdida. Esa
decisión tiene una consecuencia poco discutida: los puntajes resultantes dejan de
ser probabilidades calibradas y sobreestiman el riesgo absoluto. La corrección
estándar es una transformación monótona ajustada sobre datos no vistos, como la
regresión isotónica~\parencite{zadrozny2002}. Por ser monótona, preserva el
ordenamiento --y con él el AUROC-- mientras corrige el nivel.

\section{Equidad algorítmica}

\subsection{Definiciones}

Sea $A$ un atributo sensible, $Y$ el desenlace observado, $\hat{Y}$ la predicción
binaria y $R$ el puntaje de riesgo. Las tres familias de criterios que se emplean
en esta tesis son:

\begin{description}[leftmargin=*,style=nextline]
  \item[Paridad demográfica (independencia)]
    $P(\hat{Y}=1 \mid A=a)$ constante en $a$. Iguala la tasa de selección entre
    grupos, sin condicionar al desenlace.
  \item[\emph{Equalized odds} (separación)]
    $P(\hat{Y}=1 \mid Y=y, A=a)$ constante en $a$ para $y\in\{0,1\}$; es decir,
    igualdad simultánea de sensibilidad (TPR) y tasa de falsos positivos
    (FPR)~\parencite{hardt2016}. Su relajación a $y=1$ --igualdad de
    oportunidades-- es la que más importa clínicamente: significa que el modelo
    detecta a quien va a fallecer con la misma probabilidad en todos los grupos.
  \item[Calibración por grupo (suficiencia)]
    $P(Y=1 \mid R=r, A=a)$ constante en $a$. Significa que un riesgo de 0.3
    quiere decir lo mismo en cada entidad federativa.
\end{description}

\subsection{El teorema de imposibilidad}

Estos criterios no pueden satisfacerse a la vez. Cuando la prevalencia del
desenlace difiere entre grupos y el clasificador no es perfecto, calibración por
grupo y \emph{equalized odds} son mutuamente
incompatibles~\parencite{chouldechova2017,kleinberg2017}.

El resultado no es un tecnicismo, y conviene enunciar con precisión la forma que
toma aquí. Las tasas de mortalidad por entidad difieren por un factor de
\entFactorTasa, de modo que sobre un puntaje bien calibrado ningún conjunto de
reglas de decisión puede igualar \emph{a la vez} la sensibilidad y la tasa de
falsos positivos entre entidades. Sí puede igualarse \emph{una} de las dos: un
umbral por grupo deja intacta la calibración --que es propiedad del puntaje y no
de la regla-- y permite escoger cuál de las dos tasas se reparte por igual, a
costa de desigualar la otra. Es lo que ocurre en el capítulo~\ref{cap:resultados}.
La elección entre criterios es, por tanto, normativa y no estadística, y por eso
el entregable adecuado no es un único modelo «justo» sino una frontera que
explicite el intercambio.

\subsection{Mitigación}

Las técnicas se clasifican por la etapa del canal en la que intervienen:

\begin{description}[leftmargin=*,style=nextline]
  \item[Pre-procesamiento] Modifican los datos antes de entrenar: reponderación
    de muestras, remuestreo o corrección de representación. Son agnósticas al
    modelo, pero solo actúan sobre el desequilibrio de representación.
  \item[In-procesamiento] Incorporan la restricción de equidad en la
    optimización. La reducción por gradiente exponenciado convierte el problema
    restringido en una secuencia de problemas ponderados y admite restricciones
    de \emph{equalized odds}~\parencite{fairlearn2020}.
  \item[Post-procesamiento] Ajustan la regla de decisión sobre un modelo ya
    entrenado, típicamente eligiendo un umbral por grupo~\parencite{hardt2016}. Son
    las más eficaces cuando el problema no es el ordenamiento sino dónde se corta,
    y las más controvertidas porque usan el atributo sensible en el momento de la
    decisión.
\end{description}

Un problema abierto especialmente pertinente aquí es la equidad de subgrupos
pequeños: al desagregar por 32 entidades, los estados con menos registros
producen estimaciones de TPR con varianza alta, y los métodos que optimizan
brechas máximas tienden a perseguir ruido.

\section{Trabajos relacionados y vacío de investigación}

\subsection{La parte predictiva mexicana está saturada}

Varios trabajos revisados por pares construyen modelos de riesgo sobre la base
abierta mexicana con buen desempeño. \textcite{becerra2022} comparan modelos sobre datos mexicanos y alcanzan
alrededor de 90\,\% de exactitud, identificando neumonía, edad avanzada e
intubación como principales factores asociados a la mortalidad.
\textcite{mendez2024} aplica \emph{random forest} y XGBoost sobre
más de veinte millones de observaciones y concluye que diabetes e hipertensión,
junto con la desigualdad económica, definen la mortalidad.

La implicación metodológica es directa: replicar el esquema «clasificador más
atribución de importancia sobre comorbilidades» sería un trabajo derivativo. Por
eso el modelo predictivo se adopta aquí como línea de base y no como aportación.

\nota{La variable \texttt{INTUBADO} que aparece entre los factores más
informativos en parte de esta literatura describe el curso hospitalario y no está
disponible al primer contacto. Usarla como predictor de mortalidad en ese punto
constituye
fuga de información. Esta tesis la excluye explícitamente
(sección~\ref{sec:fuga}).}

\subsection{Lo geográfico ya se estudió, pero como predictor}

Un trabajo mexicano reciente muestra que los factores de ubicación --municipio,
unidad médica, entidad de nacimiento y de residencia-- superan incluso a la edad
como predictores de mortalidad por COVID-19 en México~\parencite{geograficos2024}. Es
un hallazgo relevante que debe reconocerse, pero responde a una pregunta
distinta.

La diferencia es sustantiva. Ese estudio usa la geografía como variable de
\emph{entrada} para responder «¿dónde vives predice tu desenlace?»,
una pregunta de epidemiología espacial. Esta tesis usa la geografía como eje de
\emph{auditoría} para responder «¿el modelo comete más errores en unas entidades
que en otras, y cómo se corrige?», una pregunta de equidad algorítmica. Que la
entidad sea informativa como predictor no dice nada sobre si el modelo es
igualmente confiable en cada entidad; son propiedades independientes.

\subsection{La equidad algorítmica en salud es madura, pero fuera de México}

La literatura de equidad en salud es rica y está anclada en otros contextos y
otros atributos sensibles. \textcite{nejadshamsi2026} auditan y
mitigan el sesgo por sexo en clasificación de severidad de COVID-19 con datos de
Canadá. En imagen médica se ha documentado sistemáticamente cómo los sesgos de
composición de las cohortes de entrenamiento se trasladan al desempeño por
subgrupo~\parencite{riccilara2022}. No se identificó trabajo alguno que use datos
mexicanos tratando la entidad federativa como atributo sensible ligado a la
desigualdad de infraestructura de salud.

\subsection{El vacío que ocupa esta tesis}

\begin{table}[htbp]
\centering
\caption{Posición de esta tesis respecto de la literatura.}
\label{tab:vacio}
\small
\begin{tabular}{p{4.6cm} p{4.6cm} p{4.6cm}}
\toprule
Dimensión & Estado en la literatura & Aporte de esta tesis \\
\midrule
Predicción de desenlace (MX) & Saturado & Línea de base, no aportación \\
Geografía como predictor (MX) & Estudiado & Se distingue: geografía como eje de equidad \\
Equidad en salud & Madura, fuera de México & Se traslada y adapta al contexto mexicano \\
Auditoría de equidad por entidad (MX) & No identificado & Eje central y original \\
Mitigación y frontera de Pareto (MX) & No identificado & Instrumento de decisión entregable \\
Validación externa de la equidad & Rara vez realizada & Pendiente (véase sección~\ref{sec:alcance}) \\
\bottomrule
\end{tabular}
\end{table}

En una frase: sabemos que la geografía predice la mortalidad en México y que la
equidad algorítmica importa en otros países, pero nadie ha preguntado todavía si
un modelo clínico entrenado con datos mexicanos es igual de preciso entre
entidades federativas, ni cómo corregir esa inequidad.
